\section{Introduction}%
\label{sec:sys:intro}

\Cref{chap:background} established that the platform and the perception pipeline
are no longer the limiting factors in warehouse inventory inspection. Multirotor
vehicles operate reliably in confined indoor spaces, and vision-based
\gls{abb:qr} decoding is accurate enough for the task. The limitation lies
elsewhere, in the assumptions these systems make before take-off: that the layout
of the warehouse is known, that it remains unchanged for the duration of the
mission, and that a single vehicle is sufficient. \Cref{chap:nav} examined the
layer at which these assumptions are made, the decision layer that determines
where a vehicle goes next, and identified four gaps. Work allocation between
several vehicles is always decided before take-off and never revised in flight.
Systems that select their next target from the meaning of the observed scene
operate as a single agent. The few systems that explore an unmapped space
autonomously have been validated only on ground robots. And those systems address
a task that terminates at the first successful detection, whereas an inventory
mission must cover the entire space.

These gaps motivate the design of a system that does not rely on a known warehouse
layout or a fixed plan before take-off. We therefore designed SwarmScan as a
cooperative system in which several vehicles build a shared representation of the
warehouse during the mission, continuously generate and select targets from this
representation, and redistribute work as the situation changes. The system also
integrates semantic guidance to help direct exploration, while its decision and
coordination mechanisms allow the mission to continue when obstacles appear,
inventory changes, or a vehicle is lost.

The chapter follows the order in which the system is constructed.
\Cref{sec:sys:problem} formulates the problem and derives the requirements a
solution must satisfy. \Cref{sec:sys:overview} presents an overview of the system
and the connections between its components.
\Cref{subsec:sys:perception,subsec:sys:map,sec:sys:candidates,sec:sys:scoring,sec:sys:claims,subsec:sys:semantic}
develop each component in turn: the two-track perception, the shared map, the
generation and scoring of targets, the coordination mechanism, and the semantic
guide. \Cref{sec:sys:mission} reassembles them and follows a complete mission
from take-off to termination. \Cref{fig:sys:architecture} summarises the
architecture developed in the following sections.

\section{Motivation and Contributions}%
\label{sec:sys:motivation}

The system designed in this chapter addresses these four gaps within a single
architecture. To the best of our knowledge, no published system combines
in-flight target redistribution, exploration of an unknown warehouse, and
semantic guidance within the same architecture.

\begin{enumerate}[itemsep=0.5em]
\item \textbf{Cooperation in an unknown warehouse}:
Several vehicles cooperate inside a warehouse whose layout is discovered
during the mission rather than supplied beforehand.

\item \textbf{Online planning}:
      No route and no division of the space are computed before take-off;
      decisions are derived continuously from the current state of the shared
      representation.

\item \textbf{A single shared representation}:
      The coverage plan normally computed on the ground is replaced by a
      single representation shared by the whole team, from which observations
      are accumulated and decisions are taken.

\item \textbf{Unified target selection}:
      Three classes of target: reading a spotted label, observing an unseen
      surface, and extending the known regio, are expressed as a common
      pose and viewing direction and scored on a single scale.

\item \textbf{In-flight target redistribution}:
      Targets are produced continuously and move between vehicles while the
      mission runs through temporary target announcements that expire unless
      renewed.

\item \textbf{Robustness to failures and changes}:
      The same mechanism handles vehicle loss, newly appearing obstacles, and
      changes in inventory by returning unfinished targets to the common list
      and updating the shared representation when new observations contradict
      earlier ones.

\item \textbf{Semantic guidance}:
      A pre-trained vision-language model is integrated into the decision loop
      and advises the team on where to search next, a role that has so far
      been evaluated only on single ground platforms.


\end{enumerate}
